\documentclass[a4paper,10pt,twoside]{article}
\usepackage{fontspec}
\setmainfont{Liberation Sans}
\setsansfont{Liberation Sans}
\renewcommand{\familydefault}{\sfdefault}
\usepackage{amsmath,amssymb}
\usepackage{geometry}
\usepackage{multicol}
\usepackage{graphicx}
\usepackage{fancyhdr}
\usepackage{lastpage}
\usepackage{booktabs}
\usepackage{listings}
\usepackage{eso-pic}
\usepackage{tikz}
\usepackage{array}
\usepackage{tabularx}
\usepackage{enumitem}
\usepackage{microtype}

\geometry{
    a4paper,
    top=1.8cm,
    bottom=1.8cm,
    left=1.2cm,
    right=1.2cm,
    headsep=14pt,
    footskip=20pt
}

\setlength{\columnsep}{16pt}
\setlength{\columnseprule}{0.35pt}

% Persistent Subtle Watermark on Every Page
\newsavebox{\watermarkbox}
\sbox{\watermarkbox}{%
    \begin{tabular}{c}
        \includegraphics[width=6.0cm]{/home/johan/GATE/GATE_CS_PYQ_DB/assets/studiva-butterfly-icon.jpg} \\[-0.2em]
        {\fontsize{14}{17}\selectfont\bfseries Studiva}
    \end{tabular}%
}
\AddToShipoutPictureBG{%
    \begin{tikzpicture}[remember picture, overlay]
        \node[opacity=0.035] at (current page.center) {\usebox{\watermarkbox}};
    \end{tikzpicture}%
}

% Running Headers & Footers
\pagestyle{fancy}
\fancyhf{}
\fancyhead[C]{\small\textbf{Studiva GATE Question Bank}}
\fancyhead[LE]{\small\textit{Computer Science \& IT (2011--2025)}}
\fancyhead[RO]{\small\textit{Studiva}}
\fancyfoot[LE,RO]{\scriptsize\textbf{Studiva} $\cdot$ \texttt{www.studiva.co.in}}
\fancyfoot[C]{\small\textbf{Page \thepage\ of \pageref{LastPage}}}
\fancyfoot[RE,LO]{\scriptsize\texttt{contact@studiva.co.in}}
\renewcommand{\headrulewidth}{0.4pt}
\renewcommand{\footrulewidth}{0.4pt}

% Listings Style (Clean, professional)
\lstset{
    basicstyle=\scriptsize\ttfamily,
    numbers=none,
    frame=single,
    rulecolor=\color{black!30},
    tabsize=2,
    breaklines=true,
    showstringspaces=false
}

\begin{document}

% ==================== PAGE 1: FORMAL TITLE & COVER PAGE ====================
\thispagestyle{empty}
\vspace*{1.5cm}
\begin{center}
    \includegraphics[width=2.5cm]{/home/johan/GATE/GATE_CS_PYQ_DB/assets/studiva-butterfly-icon.jpg}\\[1.5em]
    {\fontsize{24}{28}\selectfont\textbf{GATE COMPUTER SCIENCE \& IT}}\\[0.5em]
    {\fontsize{16}{20}\selectfont\textbf{PREVIOUS YEARS QUESTION BANK (2011 -- 2025)}}\\[0.6em]
    {\normalsize Detailed Examination Papers with Complete Official Answer Keys}\\[2.0em]
    \rule{0.65\textwidth}{0.8pt}\\[2.0em]
    {\large\textbf{Studiva GATE Examination Series}}\\[0.4em]
    {\small Comprehensive Authentic Question Bank for GATE Computer Science \& Information Technology}\\[2.8cm]

    \begin{tabular}{|ll|}
        \hline
        \textbf{Subject:} & Computer Science and Information Technology (CS/IT) \\
        \textbf{Period Covered:} & 2011 to 2025 \\
        \textbf{Examination Papers:} & 1 Complete Verified Paper Sets \\
        \textbf{Contents:} & Examination Questions \& Official Final Answer Keys \\
        \textbf{Published by:} & Studiva (\texttt{www.studiva.co.in}) \\
        \hline
    \end{tabular}

    \vfill
    {\small\textbf{Studiva $\cdot$ \texttt{www.studiva.co.in} $\cdot$ \texttt{contact@studiva.co.in}}}
\end{center}

\newpage

% ==================== PAGE 2: PUBLICATION IMPRINT & NOTES ====================
\thispagestyle{empty}
\vspace*{0.6cm}
\begin{center}
    {\large\textbf{PUBLICATION INFORMATION \& COPYRIGHT NOTICE}}\\[0.4em]
    \rule{0.6\textwidth}{0.5pt}
\end{center}

\vspace{1.2em}
\noindent\begin{tabular}{ll}
\textbf{Title:} & Studiva GATE Question Bank: Computer Science \& IT (2011--2025) \\
\textbf{Publisher:} & Studiva \\
\textbf{Website:} & \texttt{www.studiva.co.in} \\
\textbf{Official Inquiries:} & \texttt{contact@studiva.co.in} \\
\textbf{Edition:} & First Edition (2026) \\
\textbf{Syllabus Scope:} & Official GATE Syllabus (General Aptitude and CS / IT) \\
\end{tabular}

\vspace{1.8em}
\noindent\fbox{\parbox{\dimexpr\linewidth-2\fboxsep-2\fboxrule}{%
\textbf{COPYRIGHT NOTICE}\\[0.4em]
Copyright \copyright\ 2026 Studiva. All Rights Reserved.\\[0.4em]
No part of this publication may be reproduced, stored in a retrieval system, or transmitted in any form or by any means, electronic, mechanical, photocopying, recording, or otherwise, without prior written permission of Studiva, except for individual academic study.
}}

\vspace{1.5em}
\noindent\fbox{\parbox{\dimexpr\linewidth-2\fboxsep-2\fboxrule}{%
\textbf{PUBLISHER'S NOTE: ARCHIVAL QUESTION FORMAT}\\[0.4em]
This compendium has been meticulously curated and verified against official GATE question papers and final answer keys.
\par\medskip\noindent
\textbf{Notice on Image-Based Questions:}\\
Questions containing diagrams, circuit schematics, graphs, or image-based figures from the original examination papers have been omitted in this text edition due to non-readable format in historical source archives. All text-based conceptual, analytical, mathematical, algorithmic, and numerical questions are fully included with verified official keys.
}}

\vspace{1.5em}
\begin{center}
\textbf{EXAMINATION STRUCTURE \& PRACTICE GUIDELINES}
\end{center}
\begin{enumerate}[leftmargin=*]
    \item \textbf{Sectional Mark Division:} Each paper comprises General Aptitude (GA) and Computer Science (CS). Questions are clearly grouped into 1-mark and 2-mark blocks at the head of each section.
    \item \textbf{Question Types:} Includes Multiple Choice Questions (MCQ), Multiple Select Questions (MSQ), and Numerical Answer Type (NAT).
    \item \textbf{Official Answer Keys:} For self-evaluation without distraction, complete official answer keys are printed at the conclusion of each respective examination paper set.
\end{enumerate}

\newpage

% ==================== PAGE 3: TABLE OF CONTENTS ====================
\begin{center}
    {\Large\textbf{TABLE OF CONTENTS}}\\[0.4em]
    \rule{0.5\textwidth}{0.6pt}
\end{center}
\vspace{1em}

\noindent\begin{tabularx}{\linewidth}{|X|l|c|r|}
\hline
\textbf{Examination Paper} & \textbf{Session / Shift} & \textbf{Total Questions} & \textbf{Content Type} \\
\hline

GATE 2022 Computer Science & Shift CS & 61 & Questions + Answer Key \\ \hline
\end{tabularx}

\vspace{1.5em}
\noindent{\footnotesize\textit{Note: All questions in this edition have been typeset from verified digital records and cross-referenced with official answer keys released by the organising IIT/IISc institutes.}}

\newpage

% ==================== PAPER: 2022_CS ====================
\begin{center}
{\Large\textbf{GATE 2022 EXAMINATION}}\\[3pt]
{\normalsize\textbf{COMPUTER SCIENCE AND INFORMATION TECHNOLOGY}}\\[2pt]
{\small\textbf{Session / Shift: CS}}\\[5pt]
\rule{\linewidth}{0.8pt}\vspace{-4pt}
{\footnotesize\textbf{Duration: 180 Minutes} \hfill \textbf{Total Marks: 100} \hfill \textbf{Questions: 61}}\\
\vspace{-4pt}\rule{\linewidth}{0.4pt}
\end{center}
\vspace{4pt}
\noindent\fbox{\parbox{\dimexpr\linewidth-2\fboxsep-2\fboxrule}{%
\footnotesize\textbf{Examination Instructions:}\\[1pt]
1. Question Types: Multiple Choice Questions (MCQ), Multiple Select Questions (MSQ), and Numerical Answer Type (NAT).\\
2. Negative Marking Scheme: For 1-mark MCQ, $1/3$ mark is deducted; for 2-mark MCQ, $2/3$ mark is deducted. There is no negative marking for MSQ and NAT questions.\\
3. Mark Division: Questions are grouped by mark allocation within each section as specified below. Full official answer key is provided at the end of this paper.
}}
\vspace{10pt}
\begin{multicols}{2}

\vspace{8pt}
\noindent\textbf{\large\underline{GENERAL APTITUDE}}\\[2pt]
\noindent{\footnotesize\textit{[Section Block: 1-mark questions]}}
\vspace{4pt}

\noindent\textbf{Q.1} \hfill {\footnotesize\textsf{[MCQ]}}
\par\vspace{2pt}
{\bfseries The \_\_\_\_\_\_\_\_\_ is too high for it to be considered \_\_\_\_\_\_\_\_\_. \par}
\par\vspace{2pt}\noindent
\begin{tabular*}{\columnwidth}{@{\extracolsep{\fill}}p{0.46\columnwidth}p{0.46\columnwidth}@{}}
\textbf{(A)}~fair / fare & \textbf{(B)}~faer / fair \\
\addlinespace[2pt]
\textbf{(C)}~fare / fare & \textbf{(D)}~fare / fair
\end{tabular*}
\par\vspace{6pt}\hrule height 0.25pt\vspace{6pt}
\noindent\textbf{Q.2} \hfill {\footnotesize\textsf{[MCQ]}}
\par\vspace{2pt}
{\bfseries A function $y(x)$ is defined in the interval $[0, 1]$ on the $x$-axis as
\[ y(x) = \begin{cases} 2 & \text{if } 0 \le x < \frac{1}{3} \\ 3 & \text{if } \frac{1}{3} \le x < \frac{3}{4} \\ 1 & \text{if } \frac{3}{4} \le x \le 1 \end{cases} \]
Which one of the following is the area under the curve for the interval $[0, 1]$ on the $x$-axis? \par}
\par\vspace{2pt}\noindent
\begin{tabular*}{\columnwidth}{@{\extracolsep{\fill}}p{0.46\columnwidth}p{0.46\columnwidth}@{}}
\textbf{(A)}~$5/6$ & \textbf{(B)}~$6/5$ \\
\addlinespace[2pt]
\textbf{(C)}~$13/6$ & \textbf{(D)}~$6/13$
\end{tabular*}
\par\vspace{6pt}\hrule height 0.25pt\vspace{6pt}
\noindent\textbf{Q.3} \hfill {\footnotesize\textsf{[MCQ]}}
\par\vspace{2pt}
{\bfseries Let $r$ be a root of the equation $x^2 + 2x + 6 = 0$.
Then the value of the expression $(r+2)(r+3)(r+4)(r+5)$ is \par}
\par\vspace{2pt}\noindent
\begin{tabular*}{\columnwidth}{@{\extracolsep{\fill}}p{0.46\columnwidth}p{0.46\columnwidth}@{}}
\textbf{(A)}~51 & \textbf{(B)}~$-51$ \\
\addlinespace[2pt]
\textbf{(C)}~126 & \textbf{(D)}~$-126$
\end{tabular*}
\par\vspace{6pt}\hrule height 0.25pt\vspace{6pt}
\noindent\textbf{Q.4} \hfill {\footnotesize\textsf{[MCQ]}}
\par\vspace{2pt}
{\bfseries Given below are four statements:
\begin{itemize}[itemsep=0pt,topsep=2pt]
\item Statement 1: All students are inquisitive.
\item Statement 2: Some students are inquisitive.
\item Statement 3: No student is inquisitive.
\item Statement 4: Some students are not inquisitive.
\end{itemize}
From the given four statements, find the two statements that CANNOT BE TRUE simultaneously, assuming that there is at least one student in the class. \par}
\par\vspace{2pt}\noindent
\textbf{(A)}~Statement 1 and Statement 3\par\vspace{2pt}\noindent
\textbf{(B)}~Statement 1 and Statement 2\par\vspace{2pt}\noindent
\textbf{(C)}~Statement 2 and Statement 4\par\vspace{2pt}\noindent
\textbf{(D)}~Statement 3 and Statement 4
\par\vspace{6pt}\hrule height 0.25pt\vspace{6pt}

\vspace{6pt}
\noindent{\footnotesize\textit{[Section Block: 2-mark questions]}}
\vspace{3pt}

\noindent\textbf{Q.6} \hfill {\footnotesize\textsf{[MCQ]}}
\par\vspace{2pt}
{\bfseries Some people believe that ``what gets measured, improves''. Some others believe that ``what gets measured, gets gamed''. One possible reason for the difference in the beliefs is the work culture in organizations. In organizations with good work culture, metrics help improve outcomes. However, the same metrics are counterproductive in organizations with poor work culture.
Which one of the following is the CORRECT logical inference based on the information in the above passage? \par}
\par\vspace{2pt}\noindent
\textbf{(A)}~Metrics are useful in organizations with poor work culture\par\vspace{2pt}\noindent
\textbf{(B)}~Metrics are useful in organizations with good work culture\par\vspace{2pt}\noindent
\textbf{(C)}~Metrics are always counterproductive in organizations with good work culture\par\vspace{2pt}\noindent
\textbf{(D)}~Metrics are never useful in organizations with good work culture
\par\vspace{6pt}\hrule height 0.25pt\vspace{6pt}
\noindent\textbf{Q.7} \hfill {\footnotesize\textsf{[MCQ]}}
\par\vspace{2pt}
{\bfseries In a recently conducted national entrance test, boys constituted 65\% of those who appeared for the test. Girls constituted the remaining candidates and they accounted for 60\% of the qualified candidates.
Which one of the following is the correct logical inference based on the information provided in the above passage? \par}
\par\vspace{2pt}\noindent
\textbf{(A)}~Equal number of boys and girls qualified\par\vspace{2pt}\noindent
\textbf{(B)}~Equal number of boys and girls appeared for the test\par\vspace{2pt}\noindent
\textbf{(C)}~The number of boys who appeared for the test is less than the number of girls who appeared\par\vspace{2pt}\noindent
\textbf{(D)}~The number of boys who qualified the test is less than the number of girls who qualified
\par\vspace{6pt}\hrule height 0.25pt\vspace{6pt}
\noindent\textbf{Q.8} \hfill {\footnotesize\textsf{[MCQ]}}
\par\vspace{2pt}
{\bfseries A box contains five balls of same size and shape. Three of them are green coloured balls and two of them are orange coloured balls. Balls are drawn from the box one at a time. If a green ball is drawn, it is not replaced. If an orange ball is drawn, it is replaced with another orange ball.
First ball is drawn. What is the probability of getting an orange ball in the next draw? \par}
\par\vspace{2pt}\noindent
\begin{tabular*}{\columnwidth}{@{\extracolsep{\fill}}p{0.46\columnwidth}p{0.46\columnwidth}@{}}
\textbf{(A)}~$1/2$ & \textbf{(B)}~$8/25$ \\
\addlinespace[2pt]
\textbf{(C)}~$19/50$ & \textbf{(D)}~$23/50$
\end{tabular*}
\par\vspace{6pt}\hrule height 0.25pt\vspace{6pt}
\noindent\textbf{Q.9} \hfill {\footnotesize\textsf{[MCQ]}}
\par\vspace{2pt}
{\bfseries The corners and mid-points of the sides of a triangle are named using the distinct letters P, Q, R, S, T and U, but not necessarily in the same order. Consider the following statements:
\begin{itemize}[itemsep=0pt,topsep=2pt]
\item The line joining P and R is parallel to the line joining Q and S.
\item P is placed on the side opposite to the corner T.
\item S and U cannot be placed on the same side.
\end{itemize}
Which one of the following statements is correct based on the above information? \par}
\par\vspace{2pt}\noindent
\textbf{(A)}~P cannot be placed at a corner\par\vspace{2pt}\noindent
\textbf{(B)}~S cannot be placed at a corner\par\vspace{2pt}\noindent
\textbf{(C)}~U cannot be placed at a mid-point\par\vspace{2pt}\noindent
\textbf{(D)}~R cannot be placed at a corner
\par\vspace{6pt}\hrule height 0.25pt\vspace{6pt}
\noindent\textbf{Q.10} \hfill {\footnotesize\textsf{[MCQ]}}
\par\vspace{2pt}
{\bfseries A plot of land must be divided between four families. They want their individual plots to be similar in shape, not necessarily equal in area. The land has equally spaced poles, marked as dots in the figure. Two ropes, R1 and R2, are already present and cannot be moved.
What is the least number of additional straight ropes needed to create the desired plots? A single rope can pass through three poles that are aligned in a straight line. \par}
\par\vspace{2pt}\noindent
\begin{tabular*}{\columnwidth}{@{\extracolsep{\fill}}p{0.46\columnwidth}p{0.46\columnwidth}@{}}
\textbf{(A)}~2 & \textbf{(B)}~4 \\
\addlinespace[2pt]
\textbf{(C)}~5 & \textbf{(D)}~3
\end{tabular*}
\par\vspace{6pt}\hrule height 0.25pt\vspace{6pt}

\vspace{8pt}
\noindent\textbf{\large\underline{COMPUTER SCIENCE AND INFORMATION TECHNOLOGY}}\\[2pt]
\noindent{\footnotesize\textit{[Section Block: 1-mark questions]}}
\vspace{4pt}

\noindent\textbf{Q.11} \hfill {\footnotesize\textsf{[MCQ]}}
\par\vspace{2pt}
{\bfseries Which one of the following statements is TRUE for all positive functions $f(n)$? \par}
\par\vspace{2pt}\noindent
\textbf{(A)}~$f(n^2) = \Theta((f(n))^2)$, when $f(n)$ is a polynomial\par\vspace{2pt}\noindent
\textbf{(B)}~$f(n^2) = o((f(n))^2)$\par\vspace{2pt}\noindent
\textbf{(C)}~$f(n^2) = \mathcal{O}((f(n))^2)$, when $f(n)$ is an exponential function\par\vspace{2pt}\noindent
\textbf{(D)}~$f(n^2) = \Omega((f(n))^2)$
\par\vspace{6pt}\hrule height 0.25pt\vspace{6pt}
\noindent\textbf{Q.12} \hfill {\footnotesize\textsf{[MCQ]}}
\par\vspace{2pt}
{\bfseries Which one of the following regular expressions correctly represents the language of the finite automaton given below? \par}
\par\vspace{2pt}\noindent
\textbf{(A)}~$ab^*ba^* + ba^*ab^*$\par\vspace{2pt}\noindent
\textbf{(B)}~$(ab^*b + ba^*a)^*(a + b)$\par\vspace{2pt}\noindent
\textbf{(C)}~$(ab^*b + ba^*a)^*(a^*b + b^*a)$\par\vspace{2pt}\noindent
\textbf{(D)}~$(ba^*a + ab^*b)^*(ab^* + ba^*)$
\par\vspace{6pt}\hrule height 0.25pt\vspace{6pt}
\noindent\textbf{Q.13} \hfill {\footnotesize\textsf{[MCQ]}}
\par\vspace{2pt}
{\bfseries Which one of the following statements is TRUE? \par}
\par\vspace{2pt}\noindent
\textbf{(A)}~The LALR(1) parser for a grammar G cannot have reduce-reduce conflict if the LR(1) parser for G does not have reduce-reduce conflict.\par\vspace{2pt}\noindent
\textbf{(B)}~Symbol table is accessed only during the lexical analysis phase.\par\vspace{2pt}\noindent
\textbf{(C)}~Data flow analysis is necessary for run-time memory management.\par\vspace{2pt}\noindent
\textbf{(D)}~LR(1) parsing is sufficient for deterministic context-free languages.
\par\vspace{6pt}\hrule height 0.25pt\vspace{6pt}
\noindent\textbf{Q.14} \hfill {\footnotesize\textsf{[MCQ]}}
\par\vspace{2pt}
{\bfseries In a relational data model, which one of the following statements is TRUE? \par}
\par\vspace{2pt}\noindent
\textbf{(A)}~A relation with only two attributes is always in BCNF.\par\vspace{2pt}\noindent
\textbf{(B)}~If all attributes of a relation are prime attributes, then the relation is in BCNF.\par\vspace{2pt}\noindent
\textbf{(C)}~Every relation has at least one non-prime attribute.\par\vspace{2pt}\noindent
\textbf{(D)}~BCNF decompositions preserve functional dependencies.
\par\vspace{6pt}\hrule height 0.25pt\vspace{6pt}
\noindent\textbf{Q.15} \hfill {\footnotesize\textsf{[MCQ]}}
\par\vspace{2pt}
{\bfseries Consider the problem of reversing a singly linked list. Which one of the following statements is TRUE about the time complexity of algorithms that solve the above problem in $\mathcal{O}(1)$ space? \par}
\par\vspace{2pt}\noindent
\textbf{(A)}~The best algorithm for the problem takes $\Theta(n)$ time in the worst case.\par\vspace{2pt}\noindent
\textbf{(B)}~The best algorithm for the problem takes $\Theta(n \log n)$ time in the worst case.\par\vspace{2pt}\noindent
\textbf{(C)}~The best algorithm for the problem takes $\Theta(n^2)$ time in the worst case.\par\vspace{2pt}\noindent
\textbf{(D)}~It is not possible to reverse a singly linked list in $\mathcal{O}(1)$ space.
\par\vspace{6pt}\hrule height 0.25pt\vspace{6pt}
\noindent\textbf{Q.16} \hfill {\footnotesize\textsf{[MCQ]}}
\par\vspace{2pt}
{\bfseries Suppose we are given $n$ keys, $m$ hash table slots, and two simple uniform hash functions $h_1$ and $h_2$. Further suppose our hashing scheme uses $h_1$ for the odd keys and $h_2$ for the even keys. What is the expected number of keys in a slot? \par}
\par\vspace{2pt}\noindent
\begin{tabular*}{\columnwidth}{@{\extracolsep{\fill}}p{0.46\columnwidth}p{0.46\columnwidth}@{}}
\textbf{(A)}~$m/n$ & \textbf{(B)}~$n/m$ \\
\addlinespace[2pt]
\textbf{(C)}~$2n/m$ & \textbf{(D)}~$n / (2m)$
\end{tabular*}
\par\vspace{6pt}\hrule height 0.25pt\vspace{6pt}
\noindent\textbf{Q.17} \hfill {\footnotesize\textsf{[MCQ]}}
\par\vspace{2pt}
{\bfseries Which one of the following facilitates transfer of bulk data from hard disk to main memory with the highest throughput? \par}
\par\vspace{2pt}\noindent
\textbf{(A)}~DMA based I/O transfer\par\vspace{2pt}\noindent
\textbf{(B)}~Interrupt driven I/O transfer\par\vspace{2pt}\noindent
\textbf{(C)}~Polling based I/O transfer\par\vspace{2pt}\noindent
\textbf{(D)}~Programmed I/O transfer
\par\vspace{6pt}\hrule height 0.25pt\vspace{6pt}
\noindent\textbf{Q.18} \hfill {\footnotesize\textsf{[MCQ]}}
\par\vspace{2pt}
{\bfseries Let $R_1$ and $R_2$ be two 4-bit registers that store numbers in 2's complement form. For the operation $R_1 + R_2$, which one of the following values of $R_1$ and $R_2$ gives an arithmetic overflow? \par}
\par\vspace{2pt}\noindent
\textbf{(A)}~$R_1 = 1011$ and $R_2 = 1110$\par\vspace{2pt}\noindent
\textbf{(B)}~$R_1 = 1100$ and $R_2 = 1010$\par\vspace{2pt}\noindent
\textbf{(C)}~$R_1 = 0011$ and $R_2 = 0100$\par\vspace{2pt}\noindent
\textbf{(D)}~$R_1 = 1001$ and $R_2 = 1111$
\par\vspace{6pt}\hrule height 0.25pt\vspace{6pt}
\noindent\textbf{Q.19} \hfill {\footnotesize\textsf{[MCQ]}}
\par\vspace{2pt}
{\bfseries Consider the following threads, T1, T2, and T3 executing on a single processor, synchronized using three binary semaphore variables, S1, S2, and S3, operated upon using standard `wait()` and `signal()`:
\begin{center}
\begin{tabular}{p{3.5cm}p{3.5cm}p{3.5cm}}
\textbf{T1} & \textbf{T2} & \textbf{T3} \\
\texttt{while(true)\{} & \texttt{while(true)\{} & \texttt{while(true)\{} \\
\texttt{\ \ wait(S3);} & \texttt{\ \ wait(S1);} & \texttt{\ \ wait(S2);} \\
\texttt{\ \ print("C");} & \texttt{\ \ print("B");} & \texttt{\ \ print("A");} \\
\texttt{\ \ signal(S2);\}} & \texttt{\ \ signal(S3);\}} & \texttt{\ \ signal(S1);\}}
\end{tabular}
\end{center}
Which initialization of the semaphores would print the sequence BCABCABCA\dots? \par}
\par\vspace{2pt}\noindent
\textbf{(A)}~$S_1 = 1; S_2 = 1; S_3 = 1$\par\vspace{2pt}\noindent
\textbf{(B)}~$S_1 = 1; S_2 = 1; S_3 = 0$\par\vspace{2pt}\noindent
\textbf{(C)}~$S_1 = 1; S_2 = 0; S_3 = 0$\par\vspace{2pt}\noindent
\textbf{(D)}~$S_1 = 0; S_2 = 1; S_3 = 1$
\par\vspace{6pt}\hrule height 0.25pt\vspace{6pt}
\noindent\textbf{Q.20} \hfill {\footnotesize\textsf{[MCQ]}}
\par\vspace{2pt}
{\bfseries Consider the following two statements with respect to the matrices $A_{m \times n}, B_{n \times m}, C_{n \times n}$ and $D_{n \times n}$:
\begin{itemize}[itemsep=0pt,topsep=2pt]
\item Statement 1: $\text{tr}(AB) = \text{tr}(BA)$
\item Statement 2: $\text{tr}(CD) = \text{tr}(DC)$
\end{itemize}
where $\text{tr}()$ represents the trace of a matrix. Which one of the following holds? \par}
\par\vspace{2pt}\noindent
\textbf{(A)}~Statement 1 is correct and Statement 2 is wrong.\par\vspace{2pt}\noindent
\textbf{(B)}~Statement 1 is wrong and Statement 2 is correct.\par\vspace{2pt}\noindent
\textbf{(C)}~Both Statement 1 and Statement 2 are correct.\par\vspace{2pt}\noindent
\textbf{(D)}~Both Statement 1 and Statement 2 are wrong.
\par\vspace{6pt}\hrule height 0.25pt\vspace{6pt}
\noindent\textbf{Q.21} \hfill {\footnotesize\textsf{[MCQ]}}
\par\vspace{2pt}
{\bfseries What is printed by the following ANSI C program?
\begin{lstlisting}[language=C]
#include <stdio.h>
int main(int argc, char *argv[]) {
    int x = 1, z[2] = {10, 11};
    int *p = NULL;
    p = &x;
    *p = 10;
    p = &z[1];
    *(&z[0] + 1) += 3;
    printf("%d, %d, %d\n", x, z[0], z[1]);
    return 0;
}
\end{lstlisting} \par}
\par\vspace{2pt}\noindent
\begin{tabular*}{\columnwidth}{@{\extracolsep{\fill}}p{0.46\columnwidth}p{0.46\columnwidth}@{}}
\textbf{(A)}~1, 10, 11 & \textbf{(B)}~1, 10, 14 \\
\addlinespace[2pt]
\textbf{(C)}~10, 14, 11 & \textbf{(D)}~10, 10, 14
\end{tabular*}
\par\vspace{6pt}\hrule height 0.25pt\vspace{6pt}
\noindent\textbf{Q.22} \hfill {\footnotesize\textsf{[MCQ]}}
\par\vspace{2pt}
{\bfseries Consider an enterprise network with two Ethernet segments, a web server and a firewall, connected via three routers. What is the number of subnets inside the enterprise network? \par}
\par\vspace{2pt}\noindent
\begin{tabular*}{\columnwidth}{@{\extracolsep{\fill}}p{0.46\columnwidth}p{0.46\columnwidth}@{}}
\textbf{(A)}~3 & \textbf{(B)}~12 \\
\addlinespace[2pt]
\textbf{(C)}~6 & \textbf{(D)}~8
\end{tabular*}
\par\vspace{6pt}\hrule height 0.25pt\vspace{6pt}
\noindent\textbf{Q.23} \hfill {\footnotesize\textsf{[MSQ]}}
\par\vspace{2pt}
{\bfseries Which of the following statements is/are TRUE? \par}
\par\vspace{2pt}\noindent
\textbf{(A)}~Every subset of a recursively enumerable language is recursive.\par\vspace{2pt}\noindent
\textbf{(B)}~If a language $L$ and its complement $\overline{L}$ are both recursively enumerable, then $L$ must be recursive.\par\vspace{2pt}\noindent
\textbf{(C)}~Complement of a context-free language must be recursive.\par\vspace{2pt}\noindent
\textbf{(D)}~If $L_1$ and $L_2$ are regular, then $L_1 \cap L_2$ must be deterministic context-free.
\par\vspace{6pt}\hrule height 0.25pt\vspace{6pt}
\noindent\textbf{Q.24} \hfill {\footnotesize\textsf{[MSQ]}}
\par\vspace{2pt}
{\bfseries Let WB and WT be two set associative cache organizations that use LRU algorithm for cache block replacement. WB is a write back cache and WT is a write through cache. Which of the following statements is/are FALSE? \par}
\par\vspace{2pt}\noindent
\textbf{(A)}~Each cache block in WB and WT has a dirty bit.\par\vspace{2pt}\noindent
\textbf{(B)}~Every write hit in WB leads to a data transfer from cache to main memory.\par\vspace{2pt}\noindent
\textbf{(C)}~Eviction of a block from WT will not lead to data transfer from cache to main memory.\par\vspace{2pt}\noindent
\textbf{(D)}~A read miss in WB will never lead to eviction of a dirty block from WB.
\par\vspace{6pt}\hrule height 0.25pt\vspace{6pt}
\noindent\textbf{Q.25} \hfill {\footnotesize\textsf{[MSQ]}}
\par\vspace{2pt}
{\bfseries Consider the following three relations in a relational database:
\[ \text{Employee}(eId, Name), \quad \text{Brand}(bId, bName), \quad \text{Own}(eId, bId) \]
Which of the following relational algebra expressions return the set of $eIds$ who own all the brands? \par}
\par\vspace{2pt}\noindent
\textbf{(A)}~$\Pi_{eId}(\Pi_{eId, bId}(Own) / \Pi_{bId}(Brand))$\par\vspace{2pt}\noindent
\textbf{(B)}~$\Pi_{eId}(Own) - \Pi_{eId}((\Pi_{eId}(Own) \times \Pi_{bId}(Brand)) - Own)$\par\vspace{2pt}\noindent
\textbf{(C)}~$\Pi_{eId}(\Pi_{eId, bId}(Own) / \Pi_{bId}(Own))$\par\vspace{2pt}\noindent
\textbf{(D)}~$\Pi_{eId}((\Pi_{eId}(Own) \times \Pi_{bId}(Brand)) / \Pi_{bId}(Own))$
\par\vspace{6pt}\hrule height 0.25pt\vspace{6pt}
\noindent\textbf{Q.26} \hfill {\footnotesize\textsf{[MSQ]}}
\par\vspace{2pt}
{\bfseries Which of the following statements is/are TRUE with respect to deadlocks? \par}
\par\vspace{2pt}\noindent
\textbf{(A)}~Circular wait is a necessary condition for the formation of deadlock.\par\vspace{2pt}\noindent
\textbf{(B)}~In a system where each resource has more than one instance, a cycle in its wait-for graph indicates the presence of a deadlock.\par\vspace{2pt}\noindent
\textbf{(C)}~If the current allocation of resources to processes leads the system to unsafe state, then deadlock will necessarily occur.\par\vspace{2pt}\noindent
\textbf{(D)}~In the resource-allocation graph of a system, if every edge is an assignment edge, then the system is not in deadlock state.
\par\vspace{6pt}\hrule height 0.25pt\vspace{6pt}
\noindent\textbf{Q.27} \hfill {\footnotesize\textsf{[MSQ]}}
\par\vspace{2pt}
{\bfseries Which of the following statements is/are TRUE for a group $G$? \par}
\par\vspace{2pt}\noindent
\textbf{(A)}~If for all $x, y \in G$, $(xy)^2 = x^2 y^2$, then $G$ is commutative.\par\vspace{2pt}\noindent
\textbf{(B)}~If for all $x \in G$, $x^2 = 1$, then $G$ is commutative. Here, 1 is the identity element of $G$.\par\vspace{2pt}\noindent
\textbf{(C)}~If the order of $G$ is 2, then $G$ is commutative.\par\vspace{2pt}\noindent
\textbf{(D)}~If $G$ is commutative, then a subgroup of $G$ need not be commutative.
\par\vspace{6pt}\hrule height 0.25pt\vspace{6pt}
\noindent\textbf{Q.28} \hfill {\footnotesize\textsf{[NAT]}}
\par\vspace{2pt}
{\bfseries Suppose a binary search tree with 1000 distinct elements is also a complete binary tree. The tree is stored using the array representation of binary heap trees. Assuming that the array indices start with 0, the 3rd largest element of the tree is stored at index \_\_\_\_\_\_\_\_\_\_\_\_\. (Answer in integer) \par}
\par\vspace{3pt}\noindent\textbf{Answer:} \rule{3.2cm}{0.4pt}
\par\vspace{6pt}\hrule height 0.25pt\vspace{6pt}
\noindent\textbf{Q.29} \hfill {\footnotesize\textsf{[NAT]}}
\par\vspace{2pt}
{\bfseries Consider the augmented grammar with $\{+, *, (, ), id\}$ as the set of terminals:
\begin{align*}
S' &\to S \\
S &\to S + R \mid R \\
R &\to R * P \mid P \\
P &\to ( S ) \mid id
\end{align*}
If $I_0$ is the set of two LR(0) items $\{[S' \to \bullet S], [S \to \bullet S + R]\}$, then $\text{goto}(\text{closure}(I_0), +)$ contains exactly \_\_\_\_\_\_\_\_\_\_ items. (Answer in integer) \par}
\par\vspace{3pt}\noindent\textbf{Answer:} \rule{3.2cm}{0.4pt}
\par\vspace{6pt}\hrule height 0.25pt\vspace{6pt}
\noindent\textbf{Q.30} \hfill {\footnotesize\textsf{[NAT]}}
\par\vspace{2pt}
{\bfseries Consider a simple undirected graph of 10 vertices. If the graph is disconnected, then the maximum number of edges it can have is \_\_\_\_\_\_\_\_\_\_\_\. (Answer in integer) \par}
\par\vspace{3pt}\noindent\textbf{Answer:} \rule{3.2cm}{0.4pt}
\par\vspace{6pt}\hrule height 0.25pt\vspace{6pt}
\noindent\textbf{Q.31} \hfill {\footnotesize\textsf{[NAT]}}
\par\vspace{2pt}
{\bfseries Consider a relation $R(A, B, C, D, E)$ with the following three functional dependencies:
\[ AB \to C; \quad BC \to D; \quad C \to E; \]
The number of superkeys in the relation $R$ is \_\_\_\_\_\_\_\_\_\_\_\. (Answer in integer) \par}
\par\vspace{3pt}\noindent\textbf{Answer:} \rule{3.2cm}{0.4pt}
\par\vspace{6pt}\hrule height 0.25pt\vspace{6pt}
\noindent\textbf{Q.32} \hfill {\footnotesize\textsf{[NAT]}}
\par\vspace{2pt}
{\bfseries The number of arrangements of six identical balls in three identical bins is \_\_\_\_\_\_\. (Answer in integer) \par}
\par\vspace{3pt}\noindent\textbf{Answer:} \rule{3.2cm}{0.4pt}
\par\vspace{6pt}\hrule height 0.25pt\vspace{6pt}
\noindent\textbf{Q.33} \hfill {\footnotesize\textsf{[NAT]}}
\par\vspace{2pt}
{\bfseries A cache memory that has a hit rate of 0.8 has an access latency 10 ns and miss penalty 100 ns. An optimization is done on the cache to reduce the miss rate. However, the optimization results in an increase of cache access latency to 15 ns, whereas the miss penalty is not affected. The minimum hit rate (rounded off to two decimal places) needed after the optimization such that it should not increase the average memory access time is \_\_\_\_\_\_\_\_\_\_\_\. (Answer in numeric format) \par}
\par\vspace{3pt}\noindent\textbf{Answer:} \rule{3.2cm}{0.4pt}
\par\vspace{6pt}\hrule height 0.25pt\vspace{6pt}
\noindent\textbf{Q.34} \hfill {\footnotesize\textsf{[NAT]}}
\par\vspace{2pt}
{\bfseries The value of the following limit is \_\_\_\_\_\_\_\_\_\_\_\_\. (Answer in numeric format)
\[ \lim\_{x \to 0^+} \frac{x}{1 - e^{2x}} \] \par}
\par\vspace{3pt}\noindent\textbf{Answer:} \rule{3.2cm}{0.4pt}
\par\vspace{6pt}\hrule height 0.25pt\vspace{6pt}
\noindent\textbf{Q.35} \hfill {\footnotesize\textsf{[NAT]}}
\par\vspace{2pt}
{\bfseries Consider the resolution of the domain name `www.gate.org.in` by a DNS resolver. Assume that no resource records are cached anywhere across the DNS servers and that iterative query mechanism is used in the resolution. The number of DNS query-response pairs involved in completely resolving the domain name is \_\_\_\_\_\_\_\_\_\_\_\_\. (Answer in integer) \par}
\par\vspace{3pt}\noindent\textbf{Answer:} \rule{3.2cm}{0.4pt}
\par\vspace{6pt}\hrule height 0.25pt\vspace{6pt}

\vspace{6pt}
\noindent{\footnotesize\textit{[Section Block: 2-mark questions]}}
\vspace{3pt}

\noindent\textbf{Q.36} \hfill {\footnotesize\textsf{[MCQ]}}
\par\vspace{2pt}
{\bfseries Which one of the following is the closed form for the generating function of the sequence $\{a_n\}_{n \ge 0}$ defined below?
\[ a\_n = \begin{cases} n+1, & n \text{ is odd} \\ 1, & \text{otherwise} \end{cases} \] \par}
\par\vspace{2pt}\noindent
\textbf{(A)}~$\frac{x(1+x^2)}{(1-x^2)^2} + \frac{1}{1-x}$\par\vspace{2pt}\noindent
\textbf{(B)}~$\frac{x(3-x^2)}{(1-x^2)^2} + \frac{1}{1-x}$\par\vspace{2pt}\noindent
\textbf{(C)}~$\frac{2x}{(1-x^2)^2} + \frac{1}{1-x}$\par\vspace{2pt}\noindent
\textbf{(D)}~$\frac{x}{(1-x^2)^2} + \frac{1}{1-x}$
\par\vspace{6pt}\hrule height 0.25pt\vspace{6pt}
\noindent\textbf{Q.37} \hfill {\footnotesize\textsf{[MCQ]}}
\par\vspace{2pt}
{\bfseries Consider a simple undirected unweighted graph with at least three vertices. If $A$ is the adjacency matrix of the graph, then the number of 3-cycles in the graph is given by the trace of \par}
\par\vspace{2pt}\noindent
\begin{tabular*}{\columnwidth}{@{\extracolsep{\fill}}p{0.46\columnwidth}p{0.46\columnwidth}@{}}
\textbf{(A)}~$A^3$ & \textbf{(B)}~$A^3$ divided by 2 \\
\addlinespace[2pt]
\textbf{(C)}~$A^3$ divided by 3 & \textbf{(D)}~$A^3$ divided by 6
\end{tabular*}
\par\vspace{6pt}\hrule height 0.25pt\vspace{6pt}
\noindent\textbf{Q.38} \hfill {\footnotesize\textsf{[MCQ]}}
\par\vspace{2pt}
{\bfseries Which one of the following statements is FALSE? \par}
\par\vspace{2pt}\noindent
\textbf{(A)}~The TLB performs an associative search in parallel on all its valid entries using page number of incoming virtual address.\par\vspace{2pt}\noindent
\textbf{(B)}~If the virtual address of a word given by CPU has a TLB hit, but the subsequent search for the word results in a cache miss, then the word will always be present in the main memory.\par\vspace{2pt}\noindent
\textbf{(C)}~The memory access time using a given inverted page table is always same for all incoming virtual addresses.\par\vspace{2pt}\noindent
\textbf{(D)}~In a system that uses hashed page tables, if two distinct virtual addresses V1 and V2 map to the same value while hashing, then the memory access time of these addresses will not be the same.
\par\vspace{6pt}\hrule height 0.25pt\vspace{6pt}
\noindent\textbf{Q.39} \hfill {\footnotesize\textsf{[MCQ]}}
\par\vspace{2pt}
{\bfseries Let $R_i(z)$ and $W_i(z)$ denote read and write operations on a data element $z$ by a transaction $T_i$, respectively. Consider the schedule $S$ with four transactions:
\[ S: R\_4(x) R\_2(x) R\_3(x) R\_1(y) W\_1(y) W\_2(x) W\_3(y) R\_4(y) \]
Which one of the following serial schedules is conflict equivalent to $S$? \par}
\par\vspace{2pt}\noindent
\textbf{(A)}~$T_1 \to T_3 \to T_4 \to T_2$\par\vspace{2pt}\noindent
\textbf{(B)}~$T_1 \to T_4 \to T_3 \to T_2$\par\vspace{2pt}\noindent
\textbf{(C)}~$T_4 \to T_1 \to T_3 \to T_2$\par\vspace{2pt}\noindent
\textbf{(D)}~$T_3 \to T_1 \to T_4 \to T_2$
\par\vspace{6pt}\hrule height 0.25pt\vspace{6pt}
\noindent\textbf{Q.41} \hfill {\footnotesize\textsf{[MCQ]}}
\par\vspace{2pt}
{\bfseries Consider three floating point numbers A, B and C stored in registers RA, RB and RC, respectively as per IEEE-754 single precision floating point format. The 32-bit content stored in these registers (in hexadecimal form) are as follows:
\[ RA = \text{0xC1400000}, \quad RB = \text{0x42100000}, \quad RC = \text{0x41400000} \]
Which one of the following is FALSE? \par}
\par\vspace{2pt}\noindent
\begin{tabular*}{\columnwidth}{@{\extracolsep{\fill}}p{0.46\columnwidth}p{0.46\columnwidth}@{}}
\textbf{(A)}~$A + C = 0$ & \textbf{(B)}~$C = A + B$ \\
\addlinespace[2pt]
\textbf{(C)}~$B = 3C$ & \textbf{(D)}~$(B - C) > 0$
\end{tabular*}
\par\vspace{6pt}\hrule height 0.25pt\vspace{6pt}
\noindent\textbf{Q.42} \hfill {\footnotesize\textsf{[MCQ]}}
\par\vspace{2pt}
{\bfseries Consider four processes P, Q, R, and S scheduled on a CPU as per round robin algorithm with a time quantum of 4 units. The processes arrive in the order P, Q, R, S, all at time $t = 0$. There is exactly one context switch from S to Q, exactly one context switch from R to Q, and exactly two context switches from Q to R. There is no context switch from S to P. Switching to a ready process after the termination of another process is also considered a context switch. Which one of the following is NOT possible as CPU burst time (in time units) of these processes? \par}
\par\vspace{2pt}\noindent
\textbf{(A)}~P = 4, Q = 10, R = 6, S = 2\par\vspace{2pt}\noindent
\textbf{(B)}~P = 2, Q = 9, R = 5, S = 1\par\vspace{2pt}\noindent
\textbf{(C)}~P = 4, Q = 12, R = 5, S = 4\par\vspace{2pt}\noindent
\textbf{(D)}~P = 3, Q = 7, R = 7, S = 3
\par\vspace{6pt}\hrule height 0.25pt\vspace{6pt}
\noindent\textbf{Q.43} \hfill {\footnotesize\textsf{[MCQ]}}
\par\vspace{2pt}
{\bfseries What is printed by the following ANSI C program?
\begin{lstlisting}[language=C]
#include <stdio.h>
int main(int argc, char *argv[]) {
    int a[3][3][3] = {
        {1, 2, 3, 4, 5, 6, 7, 8, 9},
        {10, 11, 12, 13, 14, 15, 16, 17, 18},
        {19, 20, 21, 22, 23, 24, 25, 26, 27}
    };
    int i = 0, j = 0, k = 0;
    for(i = 0; i < 3; i++) {
        for(k = 0; k < 3; k++)
            printf("%d ", a[i][j][k]);
        printf("\n");
    }
    return 0;
}
\end{lstlisting} \par}
\par\vspace{2pt}\noindent
\textbf{(A)}~1 2 3 \\ 10 11 12 \\ 19 20 21\par\vspace{2pt}\noindent
\textbf{(B)}~1 4 7 \\ 10 13 16 \\ 19 22 25\par\vspace{2pt}\noindent
\textbf{(C)}~1 2 3 \\ 4 5 6 \\ 7 8 9\par\vspace{2pt}\noindent
\textbf{(D)}~1 2 3 \\ 13 14 15 \\ 25 26 27
\par\vspace{6pt}\hrule height 0.25pt\vspace{6pt}
\noindent\textbf{Q.44} \hfill {\footnotesize\textsf{[MCQ]}}
\par\vspace{2pt}
{\bfseries What is printed by the following ANSI C program?
\begin{lstlisting}[language=C]
#include <stdio.h>
int main(int argc, char *argv[]) {
    char a = 'P';
    char b = 'x';
    char c = (a & b) + '*';
    char d = (a | b) - '-';
    char e = (a ^ b) + '+';
    printf("%c %c %c\n", c, d, e);
    return 0;
}
\end{lstlisting} \par}
\par\vspace{2pt}\noindent
\begin{tabular*}{\columnwidth}{@{\extracolsep{\fill}}p{0.46\columnwidth}p{0.46\columnwidth}@{}}
\textbf{(A)}~z K S & \textbf{(B)}~122 75 83 \\
\addlinespace[2pt]
\textbf{(C)}~* - + & \textbf{(D)}~P x +
\end{tabular*}
\par\vspace{6pt}\hrule height 0.25pt\vspace{6pt}
\noindent\textbf{Q.45} \hfill {\footnotesize\textsf{[MCQ]}}
\par\vspace{2pt}
{\bfseries Consider solving the following system of simultaneous equations using LU decomposition:
\begin{align*}
x_1 + x_2 - x_3 &= 2 \\
2x_1 + 4x_2 - 3x_3 &= 7 \\
2x_1 + 5x_2 - 7x_3 &= 7
\end{align*}
where $L$ and $U$ are lower and upper triangular matrices. Which one of the following is the correct combination of values for $L_{32}$, $U_{33}$, and $x_1$? \par}
\par\vspace{2pt}\noindent
\textbf{(A)}~$L_{32} = 2, U_{33} = -1/2, x_1 = -1$\par\vspace{2pt}\noindent
\textbf{(B)}~$L_{32} = 2, U_{33} = 2, x_1 = -1$\par\vspace{2pt}\noindent
\textbf{(C)}~$L_{32} = -1/2, U_{33} = 2, x_1 = 0$\par\vspace{2pt}\noindent
\textbf{(D)}~$L_{32} = -1/2, U_{33} = -1/2, x_1 = 0$
\par\vspace{6pt}\hrule height 0.25pt\vspace{6pt}
\noindent\textbf{Q.46} \hfill {\footnotesize\textsf{[MSQ]}}
\par\vspace{2pt}
{\bfseries Which of the following is/are undecidable? \par}
\par\vspace{2pt}\noindent
\textbf{(A)}~Given two Turing machines $M_1$ and $M_2$, decide if $L(M_1) = L(M_2)$.\par\vspace{2pt}\noindent
\textbf{(B)}~Given a Turing machine $M$, decide if $L(M)$ is regular.\par\vspace{2pt}\noindent
\textbf{(C)}~Given a Turing machine $M$, decide if $M$ accepts all strings.\par\vspace{2pt}\noindent
\textbf{(D)}~Given a Turing machine $M$, decide if $M$ takes more than $10^{73}$ steps on every string.
\par\vspace{6pt}\hrule height 0.25pt\vspace{6pt}
\noindent\textbf{Q.47} \hfill {\footnotesize\textsf{[MSQ]}}
\par\vspace{2pt}
{\bfseries Consider the following languages:
\[ L\_1 = \{a^n w a^n \mid w \in \{a, b\}^*\} \]
\[ L\_2 = \{w x w^R \mid w, x \in \{a, b\}^*, |w|, |x| > 0\} \]
Note that $w^R$ is the reversal of the string $w$. Which of the following is/are TRUE? \par}
\par\vspace{2pt}\noindent
\textbf{(A)}~$L_1$ and $L_2$ are regular.\par\vspace{2pt}\noindent
\textbf{(B)}~$L_1$ and $L_2$ are context-free.\par\vspace{2pt}\noindent
\textbf{(C)}~$L_1$ is regular and $L_2$ is context-free.\par\vspace{2pt}\noindent
\textbf{(D)}~$L_1$ and $L_2$ are context-free but not regular.
\par\vspace{6pt}\hrule height 0.25pt\vspace{6pt}
\noindent\textbf{Q.48} \hfill {\footnotesize\textsf{[MSQ]}}
\par\vspace{2pt}
{\bfseries Consider the following languages:
\[ L\_1 = \{w w \mid w \in \{a, b\}^*\} \]
\[ L\_2 = \{a^n b^n c^m \mid m, n \ge 0\} \]
\[ L\_3 = \{a^m b^n c^n \mid m, n \ge 0\} \]
Which of the following statements is/are FALSE? \par}
\par\vspace{2pt}\noindent
\textbf{(A)}~$L_1$ is not context-free but $L_2$ and $L_3$ are deterministic context-free.\par\vspace{2pt}\noindent
\textbf{(B)}~Neither $L_1$ nor $L_2$ is context-free.\par\vspace{2pt}\noindent
\textbf{(C)}~$L_2$, $L_3$ and $L_2 \cap L_3$ all are context-free.\par\vspace{2pt}\noindent
\textbf{(D)}~Neither $L_1$ nor its complement is context-free.
\par\vspace{6pt}\hrule height 0.25pt\vspace{6pt}
\noindent\textbf{Q.49} \hfill {\footnotesize\textsf{[MSQ]}}
\par\vspace{2pt}
{\bfseries Consider a simple undirected weighted graph $G$, all of whose edge weights are distinct. Which of the following statements about the minimum spanning trees of $G$ is/are TRUE? \par}
\par\vspace{2pt}\noindent
\textbf{(A)}~The edge with the second smallest weight is always part of any minimum spanning tree of $G$.\par\vspace{2pt}\noindent
\textbf{(B)}~One or both of the edges with the third smallest and the fourth smallest weights are part of any minimum spanning tree of $G$.\par\vspace{2pt}\noindent
\textbf{(C)}~Suppose $S \subset V$ be such that $S \neq \emptyset$ and $S \neq V$. Consider the edge with the minimum weight such that one of its vertices is in $S$ and the other in $V \setminus S$. Such an edge will always be part of any minimum spanning tree of $G$.\par\vspace{2pt}\noindent
\textbf{(D)}~$G$ can have multiple minimum spanning trees.
\par\vspace{6pt}\hrule height 0.25pt\vspace{6pt}
\noindent\textbf{Q.50} \hfill {\footnotesize\textsf{[MSQ]}}
\par\vspace{2pt}
{\bfseries The following simple undirected graph is referred to as the Petersen graph. Which of the following statements is/are TRUE? \par}
\par\vspace{2pt}\noindent
\textbf{(A)}~The chromatic number of the graph is 3.\par\vspace{2pt}\noindent
\textbf{(B)}~The graph has a Hamiltonian path.\par\vspace{2pt}\noindent
\textbf{(C)}~The following graph is isomorphic to the Petersen graph.\par\vspace{2pt}\noindent
\textbf{(D)}~The size of the largest independent set of the given graph is 3.
\par\vspace{6pt}\hrule height 0.25pt\vspace{6pt}
\noindent\textbf{Q.51} \hfill {\footnotesize\textsf{[MSQ]}}
\par\vspace{2pt}
{\bfseries Consider the following recurrence:
\[ f(1) = 1; \]
\[ f(2n) = 2f(n) - 1, \quad \text{for } n \ge 1; \]
\[ f(2n+1) = 2f(n) + 1, \quad \text{for } n \ge 1. \]
Then, which of the following statements is/are TRUE? \par}
\par\vspace{2pt}\noindent
\textbf{(A)}~$f(2^n - 1) = 2^n - 1$\par\vspace{2pt}\noindent
\textbf{(B)}~$f(2^n) = 1$\par\vspace{2pt}\noindent
\textbf{(C)}~$f(5 \cdot 2^n) = 2^{n+1} + 1$\par\vspace{2pt}\noindent
\textbf{(D)}~$f(2^n + 1) = 2^n + 1$
\par\vspace{6pt}\hrule height 0.25pt\vspace{6pt}
\noindent\textbf{Q.52} \hfill {\footnotesize\textsf{[MSQ]}}
\par\vspace{2pt}
{\bfseries Which of the properties hold for the adjacency matrix $A$ of a simple undirected unweighted graph having $n$ vertices? \par}
\par\vspace{2pt}\noindent
\textbf{(A)}~The diagonal entries of $A^2$ are the degrees of the vertices of the graph.\par\vspace{2pt}\noindent
\textbf{(B)}~If the graph is connected, then none of the entries of $A^{n-1} + I_n$ can be zero.\par\vspace{2pt}\noindent
\textbf{(C)}~If the sum of all the elements of $A$ is at most $2(n - 1)$, then the graph must be acyclic.\par\vspace{2pt}\noindent
\textbf{(D)}~If there is at least a 1 in each of $A$'s rows and columns, then the graph must be connected.
\par\vspace{6pt}\hrule height 0.25pt\vspace{6pt}
\noindent\textbf{Q.53} \hfill {\footnotesize\textsf{[MSQ]}}
\par\vspace{2pt}
{\bfseries Which of the following is/are the eigenvector(s) for the matrix given below?
\[ \begin{bmatrix} -9 & 6 & -2 & 4 \\ -8 & 6 & -3 & 1 \\ -20 & 15 & -8 & 5 \\ -32 & 21 & -7 & 12 \end{bmatrix} \] \par}
\par\vspace{2pt}\noindent
\textbf{(A)}~$\begin{bmatrix} -1 \\ 1 \\ 0 \\ 1 \end{bmatrix}$\par\vspace{2pt}\noindent
\textbf{(B)}~$\begin{bmatrix} 1 \\ 0 \\ -1 \\ 0 \end{bmatrix}$\par\vspace{2pt}\noindent
\textbf{(C)}~$\begin{bmatrix} -1 \\ 0 \\ 2 \\ 2 \end{bmatrix}$\par\vspace{2pt}\noindent
\textbf{(D)}~$\begin{bmatrix} 0 \\ 1 \\ -3 \\ 0 \end{bmatrix}$
\par\vspace{6pt}\hrule height 0.25pt\vspace{6pt}
\noindent\textbf{Q.54} \hfill {\footnotesize\textsf{[MSQ]}}
\par\vspace{2pt}
{\bfseries Consider a system with 2 KB direct mapped data cache with a block size of 64 bytes. The system has a physical address space of 64 KB and a word length of 16 bits. During the execution of a program, four data words P, Q, R, and S are accessed in that order 10 times (i.e., PQRSPQRS\dots). Hence, there are 40 accesses to data cache altogether. Assume that the data cache is initially empty and no other data words are accessed by the program. The addresses of the first bytes of P, Q, R, and S are `0xA248`, `0xC28A`, `0xCA8A`, and `0xA262`, respectively. For the execution of the above program, which of the following statements is/are TRUE with respect to the data cache? \par}
\par\vspace{2pt}\noindent
\textbf{(A)}~Every access to S is a hit.\par\vspace{2pt}\noindent
\textbf{(B)}~Once P is brought to the cache it is never evicted.\par\vspace{2pt}\noindent
\textbf{(C)}~At the end of the execution only R and S reside in the cache.\par\vspace{2pt}\noindent
\textbf{(D)}~Every access to R evicts Q from the cache.
\par\vspace{6pt}\hrule height 0.25pt\vspace{6pt}
\noindent\textbf{Q.56} \hfill {\footnotesize\textsf{[NAT]}}
\par\vspace{2pt}
{\bfseries Consider the relational database with the following four schemas and their respective instances:
\[ \text{Student}(sNo, sName, dNo), \quad \text{Dept}(dNo, dName), \quad \text{Course}(cNo, cName, dNo), \quad \text{Register}(sNo, cNo) \]
\begin{lstlisting}[language=SQL]
SELECT * FROM Student AS S WHERE NOT EXISTS
    (SELECT cNo FROM Course WHERE dNo = 'D01'
     EXCEPT
     SELECT cNo FROM Register WHERE sNo = S.sNo)
\end{lstlisting}
The number of rows returned by the above SQL query is \_\_\_\_\_\_\_\_\_\_\_\. (Answer in integer) \par}
\par\vspace{3pt}\noindent\textbf{Answer:} \rule{3.2cm}{0.4pt}
\par\vspace{6pt}\hrule height 0.25pt\vspace{6pt}
\noindent\textbf{Q.58} \hfill {\footnotesize\textsf{[NAT]}}
\par\vspace{2pt}
{\bfseries Let $G = (V, E)$ be a directed graph, where $V = \{1, 2, 3, 4, 5\}$ is the set of vertices and $E$ is the set of directed edges, as defined by the adjacency matrix $A[i][j] = 1$ if $1 \le j \le i \le 5$, and 0 otherwise. A directed spanning tree of $G$ rooted at $r \in V$ is a subgraph $T$ whose undirected version is a tree and $T$ contains a directed path from $r$ to every other vertex. The number of such directed spanning trees rooted at vertex 5 is \_\_\_\_\_\_\_\_\_\_\_\_\. (Answer in integer) \par}
\par\vspace{3pt}\noindent\textbf{Answer:} \rule{3.2cm}{0.4pt}
\par\vspace{6pt}\hrule height 0.25pt\vspace{6pt}
\noindent\textbf{Q.59} \hfill {\footnotesize\textsf{[NAT]}}
\par\vspace{2pt}
{\bfseries Consider a 100 Mbps link between an earth station (sender) and a satellite (receiver) at an altitude of 2100 km. The signal propagates at a speed of $3 \times 10^8\text{ m/s}$. The time taken (in milliseconds, rounded off to two decimal places) for the receiver to completely receive a packet of 1000 bytes transmitted by the sender is \_\_\_\_\_\_\_\_\_\. (Answer in numeric format) \par}
\par\vspace{3pt}\noindent\textbf{Answer:} \rule{3.2cm}{0.4pt}
\par\vspace{6pt}\hrule height 0.25pt\vspace{6pt}
\noindent\textbf{Q.60} \hfill {\footnotesize\textsf{[NAT]}}
\par\vspace{2pt}
{\bfseries Consider the data transfer using TCP over a 1 Gbps link. Assuming that the maximum segment lifetime (MSL) is set to 60 seconds, the minimum number of bits required for the sequence number field of the TCP header, to prevent the sequence number space from wrapping around during the MSL is \_\_\_\_\_\_\_\_\_\_\_\_\. (Answer in integer) \par}
\par\vspace{3pt}\noindent\textbf{Answer:} \rule{3.2cm}{0.4pt}
\par\vspace{6pt}\hrule height 0.25pt\vspace{6pt}
\noindent\textbf{Q.61} \hfill {\footnotesize\textsf{[NAT]}}
\par\vspace{2pt}
{\bfseries A processor X1 operating at 2 GHz has a standard 5-stage RISC instruction pipeline having a base CPI (cycles per instruction) of one without any pipeline hazards. For a given program P that has 30\% branch instructions, control hazards incur 2 cycles stall for every branch. A new version of the processor X2 operating at same clock frequency has an additional branch predictor unit (BPU) that completely eliminates stalls for correctly predicted branches. There is neither any savings nor any additional stalls for wrong predictions. There are no structural hazards and data hazards for X1 and X2. If the BPU has a prediction accuracy of 80\%, the speed up (rounded off to two decimal places) obtained by X2 over X1 in executing P is \_\_\_\_\_\_\_\_\_\_\_\_\. (Answer in numeric format) \par}
\par\vspace{3pt}\noindent\textbf{Answer:} \rule{3.2cm}{0.4pt}
\par\vspace{6pt}\hrule height 0.25pt\vspace{6pt}
\noindent\textbf{Q.62} \hfill {\footnotesize\textsf{[NAT]}}
\par\vspace{2pt}
{\bfseries Consider the queues $Q_1$ containing four elements and $Q_2$ containing none. The only operations allowed on these two queues are `Enqueue(Q, element)` and `Dequeue(Q)`. The minimum number of `Enqueue` operations on $Q_1$ required to place the elements of $Q_1$ in $Q_2$ in reverse order without using any additional storage is \_\_\_\_\_\_\_\_\_\_\_\. (Answer in integer) \par}
\par\vspace{3pt}\noindent\textbf{Answer:} \rule{3.2cm}{0.4pt}
\par\vspace{6pt}\hrule height 0.25pt\vspace{6pt}
\noindent\textbf{Q.63} \hfill {\footnotesize\textsf{[NAT]}}
\par\vspace{2pt}
{\bfseries Consider two file systems A and B, that use contiguous allocation and linked allocation, respectively. A file of size 100 blocks is already stored in A and also in B. Now, consider inserting a new block in the middle of the file (between 50th and 51st block), whose data is already available in the memory. Assume that there are enough free blocks at the end of the file and that the file control blocks are already in memory. Let the number of disk accesses required to insert a block in the middle of the file in A and B be $n_A$ and $n_B$ respectively, then the value of $n_A + n_B$ is \_\_\_\_\_\_\_\_\_\. (Answer in integer) \par}
\par\vspace{3pt}\noindent\textbf{Answer:} \rule{3.2cm}{0.4pt}
\par\vspace{6pt}\hrule height 0.25pt\vspace{6pt}
\noindent\textbf{Q.64} \hfill {\footnotesize\textsf{[NAT]}}
\par\vspace{2pt}
{\bfseries Consider a demand paging system with four page frames (initially empty) and LRU page replacement policy. For the following page reference string
\[ 7, 2, 7, 3, 2, 5, 3, 4, 6, 7, 7, 1, 5, 6, 1 \]
the page fault rate, defined as the ratio of number of page faults to the number of memory accesses (rounded off to one decimal place) is \_\_\_\_\_\_\_\_\_\. (Answer in numeric format) \par}
\par\vspace{3pt}\noindent\textbf{Answer:} \rule{3.2cm}{0.4pt}
\par\vspace{6pt}\hrule height 0.25pt\vspace{6pt}
\noindent\textbf{Q.65} \hfill {\footnotesize\textsf{[NAT]}}
\par\vspace{2pt}
{\bfseries Consider the following grammar along with translation rules:
\begin{align*}
S &\to S_1 \,\#\, T & S.val &= S_1.val * T.val \\
S &\to T & S.val &= T.val \\
T &\to T_1 \,\%\, R & T.val &= T_1.val \div R.val \\
T &\to R & T.val &= R.val \\
R &\to id & R.val &= id.val
\end{align*}
Here $\#$ and $\%$ are operators and $id$ is a token that represents an integer. Using this translation scheme, the computed value of $S.val$ for root of the parse tree for the expression `20\#10\%5\#8\%2\%2` is \_\_\_\_\_\_\_\_\_\_\_\_\. (Answer in integer) \par}
\par\vspace{3pt}\noindent\textbf{Answer:} \rule{3.2cm}{0.4pt}
\par\vspace{6pt}\hrule height 0.25pt\vspace{6pt}
\end{multicols}
\vspace{8pt}
\begin{center}
\rule{0.35\linewidth}{0.5pt}\\[2pt]
\textbf{\small --- END OF QUESTION PAPER ---}\\[-8pt]
\rule{0.35\linewidth}{0.5pt}
\end{center}
\vspace{10pt}
\subsection*{OFFICIAL ANSWER KEY --- GATE 2022 (CS)}
\vspace{-4pt}
\noindent{\footnotesize\textit{Official final answer keys published by the GATE Organising Institute. For NAT questions, allowable numerical ranges are displayed where applicable. MTA = Marks To All.}}
\vspace{4pt}
\begin{center}
\small
\resizebox{\linewidth}{!}{%
\begin{tabular}{|c|c|l||c|c|l||c|c|l||c|c|l|}
\hline
\textbf{Q} & \textbf{Type} & \textbf{Key} & \textbf{Q} & \textbf{Type} & \textbf{Key} & \textbf{Q} & \textbf{Type} & \textbf{Key} & \textbf{Q} & \textbf{Type} & \textbf{Key} \\
\hline\hline
1 & MCQ & D & 18 & MCQ & B & 34 & NAT & -0.5 & 51 & MSQ & A,B,C \\ \hline
2 & MCQ & C & 19 & MCQ & C & 35 & NAT & 4 & 52 & MSQ & A \\ \hline
3 & MCQ & D & 20 & MCQ & C & 36 & MCQ & A & 53 & MSQ & A,C,D \\ \hline
4 & MCQ & A & 21 & MCQ & D & 37 & MCQ & D & 54 & MSQ & A,B,D \\ \hline
6 & MCQ & B & 22 & MCQ & C & 38 & MCQ & C & 56 & NAT & 2 \\ \hline
7 & MCQ & D & 23 & MSQ & B,C,D & 39 & MCQ & A & 58 & NAT & 24 \\ \hline
8 & MCQ & D & 24 & MSQ & A,B,D & 41 & MCQ & B & 59 & NAT & 7.07 to 7.09 \\ \hline
9 & MCQ & B & 25 & MSQ & A,B & 42 & MCQ & D & 60 & NAT & 33 \\ \hline
10 & MCQ & D & 26 & MSQ & A,D & 43 & MCQ & A & 61 & NAT & 1.42 to 1.45 \\ \hline
11 & MCQ & A & 27 & MSQ & A,B,C & 44 & MCQ & A & 62 & NAT & 0 \\ \hline
12 & MCQ & D & 28 & NAT & 509 & 45 & MCQ & D & 63 & NAT & 153 \\ \hline
13 & MCQ & D & 29 & NAT & 5 & 46 & MSQ & A,B,C & 64 & NAT & 0.6 \\ \hline
14 & MCQ & A & 30 & NAT & 36 & 47 & MSQ & A,B,C & 65 & NAT & 80 \\ \hline
15 & MCQ & A & 31 & NAT & 8 & 48 & MSQ & B,C,D &  & &  \\ \hline
16 & MCQ & B & 32 & NAT & 7 & 49 & MSQ & A,B,C &  & &  \\ \hline
17 & MCQ & A & 33 & NAT & 0.85 & 50 & MSQ & A,B,C &  & &  \\ \hline
\end{tabular}%
}
\end{center}
\end{document}